% Part: many-valued-logic
% Chapter: three-valued-logics
% Section: goedel

\documentclass[../../../include/open-logic-section]{subfiles}

\begin{document}

\olfileid{mvl}{thr}{god}

\olsection{د ګوډل منطقونه}

کرټ ګوډل د~$n$ ارزښته منطقونو يوه لړۍ وړاندې کړه۔ هر يو ئې
ټول هغه !!{formula}s رانغاړي چې په شهودي منطق کښې معتبر دي، او
پخپله د کلاسيک منطق په دننه کښې راځي۔ دلته د دې لړۍ لومړنے
د پام وړ منطق دے:

\begin{defn}\ollabel{defn:goedel}
\emph{د ګوډل $3$ ارزښته منطق}~$\LogGod$ په دې ماتريس تعريفېږي:
\begin{enumerate}
  \item معياري بياني ژبه~$\Lang L_0$ چې
  $\lfalse$, $\lnot$, $\land$, $\lor$, $\lif$.
  نښلوونکي لري۔
  \item د رښتياوالي د قيمتونو سټ~$V = \{\True, \Undef, \False\}$۔
  \item $\True$ يوازينے ټاکل شوے قيمت دے؛ يعنې~$V^+ = \{\True\}$۔
  \item د~$\lfalse$ دپاره~$\tf{\lfalse} = \False$ دے۔ د پاتې
  نښلوونکو صدق تابعې په لاندې جدولونو کښې ورکړل شوې دي:
  \begin{center}
    \begin{tabular}{c|c} 
      $\tf{\lnot}[\LogGod]$ & \\ 
      \hline  
      $\True$ & $\False$ \\ 
      $\Undef$ & $\False$ \\
      $\False$ & $\True$ 
    \end{tabular}
    \quad
    \begin{tabular}{c|ccc} 
      $\tf{\land}[\LogGod]$ & $\True$ & $\Undef$ & $\False$ \\ 
      \hline 
      $\True$ & $\True$ & $\Undef$ & $\False$ \\ 
      $\Undef$ & $\Undef$ & $\Undef$ & $\False$\\ 
      $\False$ & $\False$ & $\False$ & $\False$ 
    \end{tabular}
    \\[2ex]
    \begin{tabular}{c|ccc} 
      $\tf{\lor}[\LogGod]$ & $\True$ & $\Undef$ & $\False$ \\ 
      \hline 
      $\True$ & $\True$ & $\True$ & $\True$ \\ 
      $\Undef$ & $\True$ & $\Undef$ & $\Undef$ \\
      $\False$ & $\True$ & $\Undef$ & $\False$ 
    \end{tabular}
    \quad
    \begin{tabular}{c|ccc} 
      $\tf{\lif}[\LogGod]$ & $\True$ & $\Undef$ & $\False$ \\ 
      \hline 
      $\True$ & $\True$ & $\Undef$ & $\False$ \\ 
      $\Undef$ & $\True$ & $\True$ & $\False$  \\ 
      $\False$ & $\True$ & $\True$ & $\True$ 
    \end{tabular}
  \end{center} 
\end{enumerate}
\end{defn}

تاسو به وګورئ چې د~$\land$ او~$\lor$ صدق جدولونه د
لوکاشېوېچ او د کليني د قوي منطق له جدولونو سره يو شان دي، خو
د~$\lnot$ او~$\lif$ جدولونه په دغو منطقونو کښې سره بېل دي۔
د ګوډل په منطق کښې~$\tf{\lnot}(\Undef) = \False$ دے۔ د لوکاشېوېچ
او کليني د منطقونو پر خلاف، دلته~$\tf{\lif}(\Undef, \False) = \False$
دے؛ او د کليني د منطق پر خلاف (خو د لوکاشېوېچ د منطق په شان)،
$\tf{\lif}(\Undef,
\Undef) = \True$ دے۔

لکه چې د شهودي منطق سره پورته ياده شوې اړيکه ښيي،
$\LogGod[3]$ شهودي منطق ته نږدې دے۔ د شهودي منطق ټول معتبر
فارمولونه په~$\LogGod[3]$ کښې تاوتولوژۍ دي؛ او ډېرې کلاسيکې
تاوتولوژۍ چې په شهودي منطق کښې معتبرې نۀ دي، په~$\LogGod[3]$
کښې هم تاوتولوژۍ نۀ دي۔ د بېلګې په توګه، لاندې فارمولونه
تاوتولوژۍ نۀ دي:
\begin{align*}
  & p \lor \lnot p && (p \lif q) \lif (\lnot p \lor q) \\
  & \lnot\lnot p \lif p && \lnot(\lnot p \land \lnot q) \lif (p \lor q) \\
  & ((p \lif q) \lif p) \lif p && \lnot(p \lif q) \lif (p \land \lnot q)
\end{align*}
خو د~$\LogGod[3]$ هره تاوتولوژي په شهودي منطق کښې معتبره
نۀ ده؛ د بېلګې په توګه~$\lnot\lnot p \lor \lnot p$ يا~$(p
\lif q) \lor (q \lif p)$۔

\begin{prob}
  د صدق جدولونه ورکړئ او وښيئ چې لاندې فارمولونه د~$\LogGod[3]$
  تاوتولوژۍ دي:
  \begin{align*}
    & \lnot\lnot p \lor \lnot p\\
    & (p \lif q) \lor (q \lif p) \\
    & \lnot(p \land q) \lif (\lnot p \lor \lnot q) \\
    & (p \lif q) \lor (q \lif r) \lor (r \lif s)
  \end{align*}
\end{prob}

\begin{prob}
  د صدق جدولونه ورکړئ او وښيئ چې لاندې فارمولونه د~$\LogGod[3]$
  تاوتولوژۍ نۀ دي:
  \begin{align*}
    & (p \lif q) \lif (\lnot p \lor q) \\
    & \lnot(\lnot p \land \lnot q) \lif (p \lor q) \\
    & ((p \lif q) \lif p) \lif p \\
    & \lnot(p \lif q) \lif (p \land \lnot q)
  \end{align*}
\end{prob}

\begin{prob}
  له لاندې اړيکو کومې د ګوډل په منطق کښې صدق کوي؟ د هرې يوې
  دپاره د صدق جدول ورکړئ۔
  \begin{enumerate}
    \item $p, p \lif q \Entails q$
    \item $p \lor q, \lnot p \Entails q$
    \item $p \land q \Entails p$
    \item $p \Entails p \land p$
    \item $p \Entails p \lor q$
  \end{enumerate}
\end{prob}

\end{document}
